% STATUS: rewritten 2026-11-05; every claim is cited.
\section{Gamma-rings}\label{sec:other-systems:gamma-rings}

Recall that a near-ring keeps only one distributive law.
A $\Gamma$-ring goes further: it replaces the product by a family of
products \cite[p.~2]{DeyPaulRakhimov2012}. There is one product for each
element of a set $\Gamma$ \cite[p.~2]{DeyPaulRakhimov2012}. Nobusawa
introduced the
notion in 1964; Dey, Paul and Rakhimov state it in the current form
\cite[p.~2]{DeyPaulRakhimov2012}.

Every derivation below is taken with respect to the whole family of products, so the rule must hold for each element of $\Gamma$ separately.

\begin{definition}[{\cite[p.~2]{DeyPaulRakhimov2012}}]\label{def:other:gamma-ring}
Let $M$ and $\Gamma$ be additive abelian groups. Then $M$ is a
\emph{$\Gamma$-ring} if there is a mapping
$M \times \Gamma \times M \to M$, written $(a, \alpha, b) \mapsto a\,\alpha\,b$,
such that for all $a, b, c \in M$ and $\alpha, \beta \in \Gamma$:
\textup{(i)} $a\,\alpha\,b \in M$; \textup{(ii)}
$(a+b)\,\alpha\,c = a\,\alpha\,c + b\,\alpha\,c$,
$a\,(\alpha+\beta)\,b = a\,\alpha\,b + a\,\beta\,b$,
$a\,\alpha\,(b+c) = a\,\alpha\,b + a\,\alpha\,c$, and
$(a\,\alpha\,b)\,\beta\,c = a\,\alpha\,(b\,\beta\,c)$.
\end{definition}

The next definition states the derivations on these rings.

\begin{definition}[{\cite[pp.~2--3]{DeyPaulRakhimov2012}}]\label{def:other:gamma-derivation}
Let $M$ be a $\Gamma$-ring. An additive mapping $D \colon M \to M$ is a
\emph{derivation} of $M$ if
\[
D(x\,\alpha\,y) = D(x)\,\alpha\,y + x\,\alpha\,D(y)
\qquad (x, y \in M,\ \alpha \in \Gamma).
\]
It is a \emph{Jordan derivation} if
$D(x\,\alpha\,x) = D(x)\,\alpha\,x + x\,\alpha\,D(x)$ for all $x \in M$
and $\alpha \in \Gamma$.
\end{definition}

The next definition allows a second derivation on one side.

\begin{definition}[{\cite[p.~3]{DeyPaulRakhimov2012}}]\label{def:other:gamma-generalized}
An additive mapping $D \colon M \to M$ is a \emph{generalized derivation}
of $M$ if there is a derivation $d$ of $M$ such that
\[
D(x\,\alpha\,y) = D(x)\,\alpha\,y + x\,\alpha\,d(y)
\qquad (x, y \in M,\ \alpha \in \Gamma).
\]
It is a \emph{Jordan generalized derivation} if
$D(x\,\alpha\,x) = D(x)\,\alpha\,x + x\,\alpha\,d(x)$ for all $x \in M$
and $\alpha \in \Gamma$.
\end{definition}

The main theorem of Dey, Paul and Rakhimov says that every Jordan
generalized derivation is generalized, under a restriction on $M$
\cite[Abstract, p.~1]{DeyPaulRakhimov2012}. The
restriction is the assumption
\begin{equation}\label{eq:other:star}
a\,\alpha\,b\,\beta\,c = a\,\beta\,b\,\alpha\,c
\qquad (a, b, c \in M,\ \alpha, \beta \in \Gamma),
\end{equation}
which the authors impose on the $\Gamma$-ring $M$ (Theorem~\ref{thm:other:gamma-main})
\cite[pp.~1--2]{DeyPaulRakhimov2012}.\begin{theorem}[{\cite[Abstract, p.~1]{DeyPaulRakhimov2012}}]\label{thm:other:gamma-main}
Let $M$ be a $2$-torsion-free semiprime $\Gamma$-ring satisfying
\eqref{eq:other:star}. Let $D \colon M \to M$ be an additive mapping such
that, for some derivation $d$ of $M$,
$D(x\,\alpha\,x) = D(x)\,\alpha\,x + x\,\alpha\,d(x)$ for all $x \in M$
and $\alpha \in \Gamma$. Then $D$ is a generalized derivation of $M$.
\end{theorem}

The theorem lifts the same statement for rings, due to Hvala, to the
$\Gamma$ setting \cite[p.~1]{DeyPaulRakhimov2012}. The next section moves
to semirings, where a different obstacle appears: addition has no
negatives at all \cite[p.~625]{AhmedDudek2021}.

A prime $\Gamma$-near-ring is shown to be a commutative $\Gamma$-ring under conditions imposed by a symmetric biderivation~\cite{WlOnCommutativityIn2014}.

A left derivation $d$ of a $\Gamma$-ring $M$ with $d(U)\subset U$ and $d^2(U)\subset Z(M)$ on an ideal $U$ is shown to force $M$ to be commutative~\cite{WlTheCommutativityInPrimeGammaRi2018}.

Chakraborty and Paul prove that every Jordan generalized $k$-derivation of a 2-torsion free prime gamma ring in the sense of Nobusawa is a generalized $k$-derivation, the converse direction being immediate~\cite{WlOnJordanGeneralizedKDerivation2007}.

U\u{g}kun, \"Ozt\"urk and Jun give conditions under which a $\Gamma$-near-ring carrying a derivation is commutative~\cite{WlOnPrimeGammaNearRingsWithDeriv2004}.

Dey, Paul and Rakhimov investigate when a prime $\Gamma$-near-ring is commutative by means of symmetric bi-derivations, and prove that under stated conditions a prime $\Gamma$-near-ring is a commutative $\Gamma$-ring~\cite{WlOnCommutativityInPrimeGammaNea2014}.
